% =====================================================================
%  Figure: the cluster (immigration--birth) representation (LEAD).
%  Schematic -- event times are hand-placed, NOT data or a simulation.
%  Panel (b) is computed exactly in TikZ from the exponential-kernel
%  intensity, using the Markov recursion between event times.
%  Requires the house colours defined in preamble.tex.
% =====================================================================
\begin{figure}[t]
\centering
\begin{tikzpicture}[x=0.952cm,y=1cm,
    ev/.style={circle,inner sep=0pt,minimum size=5.2pt},
    bg/.style={ev,fill=LabNavyLight,draw=white,line width=0.5pt},
    sh/.style={ev,fill=LabAccent,draw=white,line width=0.5pt},
    of/.style={ev,fill=LabNavy,draw=white,line width=0.5pt},
    kid/.style={LabNavyMid!75,line width=0.55pt,
                -{Stealth[length=3.2pt,width=2.6pt]},shorten >=1.6pt,shorten <=1.4pt},
    genlab/.style={font=\scriptsize,text=LabGrey,anchor=east,inner sep=1pt},
    panel/.style={font=\footnotesize\bfseries,text=LabNavy,anchor=south west,inner sep=0pt},
    note/.style={font=\scriptsize\itshape,text=LabGrey,inner sep=1pt},
    leg/.style={font=\scriptsize,text=LabInk,anchor=west,inner sep=1pt},
    axis/.style={LabNavy,line width=0.5pt},
  ]
  % ---------------- parameters (schematic) ----------------
  \def\gy{0.68}       % vertical spacing of generations in panel (a)
  \def\tshock{5.6}    % time of the salient external event
  \def\tend{15.3}     % right end of the time axis

  % =========================== panel (a) ===========================
  \node[panel] at (-1.05,3*\gy+0.55) {(a)\enspace Who triggered whom
    \textnormal{\textcolor{LabGrey}{(latent, unobserved)}}};
  \foreach \g in {0,1,2,3}{
    \draw[LabRule,line width=0.3pt,dash pattern=on 1pt off 1.5pt] (0,\g*\gy) -- (\tend,\g*\gy);
    \node[genlab] at (-0.12,\g*\gy) {gen.\,\g};}

  % a highlighted cluster (one immigrant and all its descendants)
  \begin{scope}[on background layer]
    \fill[LabTint,rounded corners=4pt] (0.38,-0.2) rectangle (2.55,2*\gy+0.22);
  \end{scope}
  \node[note,anchor=south] at (1.47,2*\gy+0.22) {one cluster};

  % background immigrants (generation 0)
  \node[bg] (b1) at (0.6,0) {};
  \node[bg] (b2) at (3.4,0) {};
  \node[bg] (b3) at (10.6,0) {};
  \node[bg] (b4) at (13.6,0) {};
  % shock-triggered immigrants (generation 0)
  \node[sh] (s1) at (5.75,0) {};
  \node[sh] (s2) at (6.05,0) {};
  \node[sh] (s3) at (6.4,0) {};
  \node[sh] (s4) at (6.95,0) {};
  \node[sh] (s5) at (7.7,0) {};
  % offspring
  \node[of] (o1) at (1.3,\gy) {};
  \node[of] (o2) at (2.2,2*\gy) {};
  \node[of] (o3) at (4.3,\gy) {};
  \node[of] (o4) at (14.5,\gy) {};
  \node[of] (c1) at (6.5,\gy) {};
  \node[of] (c2) at (7.3,\gy) {};
  \node[of] (c3) at (6.85,\gy) {};
  \node[of] (c4) at (7.05,2*\gy) {};
  \node[of] (c5) at (8.3,2*\gy) {};
  \node[of] (c6) at (8.9,3*\gy) {};
  \node[of] (c9) at (8.6,2*\gy) {};
  \node[of] (c8) at (9.4,\gy) {};
  \node[of] (c10) at (10.1,2*\gy) {};
  \node[of] (c11) at (11.4,\gy) {};
  \node[of] (c12) at (12.3,2*\gy) {};
  % parent -> child
  \foreach \p/\c in {b1/o1,o1/o2,b2/o3,b4/o4,s1/c1,s1/c2,s2/c3,c1/c4,c1/c5,
                     c4/c6,c2/c9,s4/c8,c8/c10,s5/c11,c11/c12}{
    \draw[kid] (\p) to[out=72,in=208,looseness=0.85] (\c);}

  % legend (one row, right-aligned with the time axis)
  \begin{scope}[node distance=1.5pt]
    \node[leg,anchor=east] (l3) at (\tend+0.25,3*\gy+0.66) {offspring};
    \node[of,left=of l3] (d3) {};
    \node[leg,anchor=east] (l2) at ([xshift=-9pt]d3.west) {shock-triggered immigrant};
    \node[sh,left=of l2] (d2) {};
    \node[leg,anchor=east] (l1) at ([xshift=-9pt]d2.west) {background immigrant};
    \node[bg,left=of l1] {};
  \end{scope}

  % =========================== panel (b) ===========================
  \begin{scope}[yshift=-3.75cm,y=0.45cm]
    \node[panel] at (-1.05,5.6) {(b)\enspace What is recorded: event times, and
      the conditional intensity $\lambda(t)$ they imply under the model};
    % model parameters (schematic): background, kernel, shock response
    \def\kmu{0.6}\def\kalpha{0.55}\def\kbeta{1.0}\def\kgam{0.8}\def\kjump{4.0}
    \xdef\tp{0}\xdef\Hs{0}\xdef\Ss{0}
    % breakpoints: time/kind, kind 1 = incident, 2 = shock, 0 = end of axis
    \foreach \t/\k in {0.6/1,1.3/1,2.2/1,3.4/1,4.3/1,5.6/2,5.75/1,6.05/1,6.4/1,
                       6.5/1,6.85/1,6.95/1,7.05/1,7.3/1,7.7/1,8.3/1,8.6/1,8.9/1,
                       9.4/1,10.1/1,10.6/1,11.4/1,12.3/1,13.6/1,14.5/1,\tend/0}{
      \begin{scope}[on background layer]
        \filldraw[LabNavyLight!45,line width=0.25pt] (\tp,0) rectangle (\t,\kmu);
        \filldraw[LabAccent!55,line width=0.25pt]
          plot[domain=\tp:\t,samples=24] (\x,{\kmu+\Ss*exp(-\kgam*(\x-\tp))})
          -- (\t,\kmu) -- (\tp,\kmu) -- cycle;
        \filldraw[LabNavyMid!80,line width=0.25pt]
          plot[domain=\tp:\t,samples=24]
            (\x,{\kmu+\Ss*exp(-\kgam*(\x-\tp))+\Hs*exp(-\kbeta*(\x-\tp))})
          -- plot[domain=\t:\tp,samples=24] (\x,{\kmu+\Ss*exp(-\kgam*(\x-\tp))})
          -- cycle;
      \end{scope}
      \draw[LabNavy,line width=0.7pt,line join=round]
        plot[domain=\tp:\t,samples=24]
          (\x,{\kmu+\Ss*exp(-\kgam*(\x-\tp))+\Hs*exp(-\kbeta*(\x-\tp))});
      \pgfmathsetmacro{\Hn}{\Hs*exp(-\kbeta*(\t-\tp))+(\k==1)*\kalpha*\kbeta}
      \pgfmathsetmacro{\Sn}{\Ss*exp(-\kgam*(\t-\tp))+(\k==2)*\kjump}
      \pgfmathsetmacro{\Lold}{\kmu+\Ss*exp(-\kgam*(\t-\tp))+\Hs*exp(-\kbeta*(\t-\tp))}
      \pgfmathsetmacro{\Lnew}{\kmu+\Sn+\Hn}
      \ifnum\k>0 \draw[LabNavy,line width=0.7pt] (\t,\Lold) -- (\t,\Lnew);\fi
      \ifnum\k=1 \draw[LabInk,line width=0.6pt] (\t,-0.12) -- (\t,-0.85);\fi
      \xdef\tp{\t}\xdef\Hs{\Hn}\xdef\Ss{\Sn}
    }
    % axes
    \draw[axis,-{Stealth[length=4pt]}] (0,0) -- (\tend+0.25,0)
      node[right,font=\footnotesize,text=LabNavy,inner sep=1pt] {$t$};
    \draw[axis] (0,0) -- (0,5.2);
    \node[genlab,anchor=east] at (-0.1,4.6) {$\lambda(t)$};
    \node[genlab,anchor=east] at (-0.1,\kmu) {$\mu$};
    \node[genlab,anchor=east] at (-0.1,-0.5) {events};
    % layer labels (direct labelling, short leaders)
    \node[font=\scriptsize,text=LabAccent,anchor=east,inner sep=1pt] (lshock)
      at (5.3,3.3) {shock response $\beta_kS_e\,\psi_k(t-\tau_e)$};
    \draw[LabAccent,line width=0.4pt] (lshock.east) -- (6.05,2.0);
    \fill[LabAccent] (6.05,2.0) circle (0.9pt);
    \node[font=\scriptsize,text=LabNavyMid,anchor=west,inner sep=1pt] (lexc)
      at (10.3,3.6) {excitation by earlier incidents $\textstyle\sum_{t_i<t}\varphi(t-t_i)$};
    \draw[LabNavyMid,line width=0.4pt] (lexc.west) -- (9.45,1.45);
    \fill[white] (9.45,1.45) circle (0.9pt);
    \node[font=\scriptsize,text=LabNavy,anchor=west,inner sep=1pt] at (12.35,0.3) {background $\mu$};
  \end{scope}

  % the salient external event, spanning both panels
  \draw[LabAccent,line width=0.7pt,dash pattern=on 2.5pt off 1.5pt]
    (\tshock,-3.75-0.45) -- (\tshock,-3.75+5.0*0.45)
    (\tshock,-0.32) -- (\tshock,3*\gy+0.12);
  \node[font=\scriptsize,text=LabAccent,anchor=south,align=center,inner sep=1pt]
    at (\tshock,3*\gy+0.12) {salient external event $\tau_e$};
\end{tikzpicture}
\caption{\textbf{The cluster representation of a Hawkes process with
exogenous shocks.} \emph{(a)} Every incident is either an \emph{immigrant}
or the \emph{offspring} of an earlier incident. Background immigrants (light
navy) arrive at rate $\mu_k(t)$; shock-triggered immigrants (orange) are the
direct response to a salient external event at $\tau_e$; each incident, of
any kind, independently produces $\Pois(\alpha_{kj})$ offspring at delays
drawn from $g_{kj}$ (navy), which produce offspring in turn
\citep{hawkes1974}. An immigrant and all its descendants form one cluster.
\emph{(b)} The data are the event times alone (ticks). Under the model, the
conditional intensity is the sum of the background, the shock response and
the excitation by past incidents (\cref{eq:th-intensity}). The share of each
layer at an event time is the probability that the event is a background
immigrant, a shock-triggered immigrant or the offspring of an earlier
incident; this is how the EM algorithm reconstructs the unobserved family
trees (\cref{alg:em}). The surge decomposition of \cref{sec:attribution}
classifies incidents instead by the root of their cluster
(\cref{fig:attribution}). Schematic: event times are
hand-placed, and the intensity is computed exactly from them for
illustrative parameters ($\mu=0.6$, $\alpha=0.55$, exponential kernel
with $\beta=1$, and $\beta_kS_e=5$ expected direct incidents with
response $\psi_k(t)=\eta e^{-\eta t}$, $\eta=0.8$).}
\label{fig:branching}
\end{figure}
